现在设 $G$ 是任意群, $M$ 是任意集合. 若存在从 $G$ 到 $M$ 到自身的双射映射所成的群的同态 $\phi$, 就称 \phantomsection\label{idx:group-action-definition}\textbf{$G$ 作用于 $M$}. 等价地, 写作
\[
  \phi(g)(p)=A(g,p),
\]
则要求存在满足以下条件的映射 $A:G\times M\to M$:
\begin{enumerate}[label=(\roman*),leftmargin=*]
\item 若 $e$ 是 $G$ 的单位元, 则 $A(e,p)=p$, $\forall p\in M$;
\item 若 $g,h\in G$, 则 $A(g,A(h,p))=A(gh,p)$, $\forall p\in M$.
\end{enumerate}
通常用 $g\cdot p$ 表示 $A(g,p)$.

\begin{RiemannExample}\label{ex:group-actions-by-left-multiplication}
\begin{enumerate}[label=(\arabic*),leftmargin=*]
\item 设 $G$ 是群, $H\subset G$ 是子群. 则 $H$ 通过左乘作用于 $G$: 对 $h\in H$, $g\in G$, 有 $A(h,g)=h\cdot g$.
\item $GL(n)$ 通过 $A\cdot x=Ax$ 作用于 $\mathbb{R}^n$, 其中 $A\in GL(n)$ 且 $x\in\mathbb{R}^n$. 对任何子群 $G\subset GL(n)$ 也是如此.
\end{enumerate}
\end{RiemannExample}

对每个 $p\in M$, 可以把 $p$ 的\textbf{轨道}定义为集合 $G\cdot p:=\{g\cdot p\mid g\in G\}$. 若 $G\cdot p=\{p\}$, 则称 $p$ 为 $G$ 的\textbf{不动点}. 若存在点 $p\in M$ 的轨道是整个 $M$ (即 $G\cdot p=M$), 则称作用是\textbf{传递的}. 注意, 此时只有一条轨道, 且对任意 $p,q\in M$ 且 $p\ne q$, 总存在群的元素 $g\in G$ 使得 $q=g\cdot p$. 这时称流形 $M$ 为 $G$ 的\textbf{齐性空间}. 点 $p\in M$ 的\textbf{稳定子} (或\textbf{各向同性子群}) 是群
\[
  G_p=\{g\in G\mid g\cdot p=p\}.
\]
若所有稳定子都是平凡的, 则称作用是\textbf{自由的}.

若 $G$ 是 Lie 群, $M$ 是光滑流形, 当映射 $A:G\times M\to M$ 可微时, 称作用是\textbf{光滑的}. 此时, 映射 $p\mapsto g\cdot p$ 是微分同胚. 我们始终假设 Lie 群在可微流形上的作用是光滑的. 若映射
\begin{align*}
  G\times M&\to M\times M\\
  (g,p)&\mapsto(g\cdot p,p)
\end{align*}
是固有映射, 则称这个光滑作用是\textbf{固有的} (回顾一下, 若每个紧集的原像都紧, 则称映射是\textbf{固有映射}, 参见第~\RGref{subsec:notes-proper-maps-and-fundamental-group}~节).

\begin{RiemannRemark}\label{rem:sequential-characterization-of-proper-actions}
注意, 一个光滑作用是固有的, 当且仅当给定 $M$ 中两个收敛序列 $\{p_n\}$ 和 $\{g_n\cdot p_n\}$, $G$ 中存在收敛子序列 $\{g_{n_k}\}$. 若 $G$ 紧致, 这个条件总是成立.
\end{RiemannRemark}

$G$ 在 $M$ 上的作用的轨道是由 $p\sim q\Leftrightarrow q\in G\cdot p$ 给出的等价关系 $\sim$ 的等价类 (参见第~\RGref{subsec:notes-topological-manifolds}~节). 因此, 商 (拓扑) 空间 $M/{\sim}$ 通常称为该作用的\textbf{轨道空间}, 并记作 $M/G$.

\begin{RiemannProposition}\label{prop:hausdorff-orbit-space}
若 Lie 群 $G$ 在可微流形 $M$ 上的作用是固有的, 则轨道空间 $M/G$ 是 Hausdorff 空间.
\end{RiemannProposition}

\begin{RiemannProof}
关系 $p\sim q\Leftrightarrow q\in G\cdot p$ 是开等价关系 (参见第~\RGref{subsec:notes-topological-manifolds}~节). 实际上, 由于 $p\mapsto g\cdot p$ 是同胚, 对 $M$ 中的任何开集 $U$, 集合 $\pi^{-1}(\pi(U))=\{g\cdot p\mid p\in U\text{ 且 }g\in G\}=\bigcup_{g\in G}g\cdot U$ 是 $M$ 的开子集, 即 $\pi(U)$ 是开集 (这里 $\pi:M\to M/G$ 是商映射). 因此只需证明集合
\[
  R=\{(p,q)\in M\times M\mid p\sim q\}
\]
是闭的 (参见 Proposition~\RGref{prop:topology-of-quotients}). 这一点由 $R$ 是映射
\begin{align*}
  G\times M&\to M\times M\\
  (g,p)&\mapsto(g\cdot p,p)
\end{align*}
的像推出, 因为该映射连续而且固有, 因而是闭映射 (参见第~\RGref{subsec:notes-proper-maps-and-fundamental-group}~节).
\end{RiemannProof}

在一定条件下, 轨道空间 $M/G$ 自然成为可微流形.

\begin{RiemannTheorem}\label{thm:quotient-manifold-of-free-proper-action}
设 $M$ 是可微流形, 赋有 Lie 群 $G$ 的一个自由且固有的作用. 则轨道空间 $M/G$ 自然成为维数为 $\dim M-\dim G$ 的可微流形, 而且商映射 $\pi:M\to M/G$ 是淹没.
\end{RiemannTheorem}

\begin{RiemannProof}
由前一命题, 商空间 $M/G$ 是 Hausdorff 的. 此外, 由于 $M$ 满足第二可数公理, 且 $G$ 所定义的等价关系是开等价关系, 这个商空间也满足第二可数公理. 还需证明 $M/G$ 具有一种自然的可微结构, 使商映射是淹没. 我们只证明 $G$ 是离散 (即零维) Lie 群的情形 (参见 Remark~\ref{rem:topological-manifolds-low-dimension}); 一般 Lie 群的证明可见 \RGcite{DK99}.

在这里, 只需证明: 对每个 $p\in M$, 存在邻域 $U\ni p$, 使得当 $g\ne h$ 时, $g\cdot U\cap h\cdot U=\varnothing$. 这保证每个点 $[p]\in M/G$ 都有一个与 $U$ 同胚的邻域 $[U]$, 并可假定 $U$ 是坐标邻域. 由于 $G$ 通过微分同胚作用, 如此定义的可微结构与 $p\in[p]$ 的选取无关. $M/G$ 的坐标图来自 $M$ 的坐标图, 所以重叠映射是光滑的. 因此, $M/G$ 具有一种自然的可微结构, 使 $\pi:M\to M/G$ 是局部微分同胚 (因为 $\pi|_U:U\to[U]$ 的坐标表达式是恒等映射).

证明当 $g\ne h$ 时 $g\cdot U\cap h\cdot U=\varnothing$, 等价于证明当 $g\ne e$ 时 $g\cdot U\cap U=\varnothing$. 假设对于任何邻域 $U\ni p$, 这一点都不成立. 则存在一列开集 $U_n\ni p$, 满足 $U_{n+1}\subset U_n$, $\bigcap_{n=1}^{+\infty}U_n=\{p\}$, 以及一列 $g_n\in G\setminus\{e\}$, 使得 $g_n\cdot U_n\cap U_n\ne\varnothing$. 取 $p_n\in g_n\cdot U_n\cap U_n$. 则对某个 $q_n\in U_n$, 有 $p_n=g_n\cdot q_n$. 我们有 $p_n\to p$ 及 $q_n\to p$. 由于作用固有, $g_n$ 存在收敛子序列 $g_{n_k}$. 设 $g$ 为其极限. 在 $q_{n_k}=g_{n_k}\cdot p_{n_k}$ 中令 $k\to+\infty$, 得 $g\cdot p=p$, 从而 $g=e$ (因为作用是自由的). 由于 $G$ 离散, 对充分大的 $k$ 应有 $g_{n_k}=e$, 矛盾.
\end{RiemannProof}

\begin{RiemannExample}\label{ex:projective-spaces-and-tori-as-quotients}
\begin{enumerate}[label=(\arabic*),ref=\arabic*,leftmargin=*]
\item\label{ex:projective-space-antipodal-quotient} 设 $S^n=\{x\in\mathbb{R}^{n+1}\mid\sum_{i=1}^{n}(x^i)^2=1\}$ 赋有 $G=\mathbb{Z}_2=\{-I,I\}$ 由 $-I\cdot x=-x$ (对径映射) 给出的作用. 这个作用既固有又自由, 所以轨道空间 $S^n/G$ 是 $n$ 维流形. 这个空间就是实射影空间 $\mathbb{R}P^n$ [参见 Exercise~\ref{ex:exercise-differentiable-structures}(\ref{ex:projective-space-atlas-from-sphere})].
\item 群 $G=\mathbb{R}\setminus\{0\}$ 通过乘法 $t\cdot x=tx$ 作用于 $M=\mathbb{R}^{n+1}\setminus\{0\}$. 这个作用既固有又自由, 因而 $M/G$ 是 $n$ 维可微流形 (它又是 $\mathbb{R}P^n$).
\item 考虑 $M=\mathbb{R}^n$, 赋有 $G=\mathbb{Z}^n$ 由下式定义的作用:
\[
  (k^1,\ldots,k^n)\cdot(x^1,\ldots,x^n)
  =(x^1+k^1,\ldots,x^n+k^n).
\]
这个作用既固有又自由, 所以商空间 $M/G$ 是 $n$ 维流形. 赋予 Theorem~\ref{thm:quotient-manifold-of-free-proper-action} 所定义的商可微结构后, 这个空间称为 \textbf{$n$ 维环面}, 记为 $T^n$. 它与乘积流形 $S^1\times\cdots\times S^1$ 微分同胚, 且当 $n=2$ 时, 它与 $\mathbb{R}^3$ 中的旋转环面微分同胚.
\end{enumerate}
\end{RiemannExample}

离散群作用的商给出了流形的\textbf{覆盖}.

\begin{RiemannDefinition}\label{def:smooth-covering-of-manifold}
可微流形 $B$ 的一个\textbf{光滑覆盖}是二元组 $(M,\pi)$, 其中 $M$ 是连通可微流形, $\pi:M\to B$ 是满射的局部微分同胚, 且对每个 $p\in B$, 在 $B$ 中都存在 $p$ 的连通邻域 $U$, 使得 $\pi^{-1}(U)$ 是互不相交的开集 $U_\alpha\subset M$ (称为\textbf{片}) 的并, 且 $\pi$ 到 $U_\alpha$ 的限制 $\pi_\alpha$ 都是到 $U$ 的微分同胚. 映射 $\pi$ 称为\textbf{覆盖映射}, $M$ 称为\textbf{覆盖流形}.
\end{RiemannDefinition}

\begin{RiemannRemark}\label{rem:dimensions-and-fibers-of-coverings}
\begin{enumerate}[label=(\arabic*),leftmargin=*]
\item 显然, 必有 $\dim M=\dim B$.
\item 注意, 互不相交的开集族 $\{U_\alpha\}$ 必为可数 ($M$ 有可数基).
\item \textbf{纤维} $\pi^{-1}(p)\subset M$ 具有离散拓扑. 实际上, 每一片 $U_\alpha$ 都是开集且与 $\pi^{-1}(p)$ 恰好相交于一点, 因此该点在子空间拓扑中是开集.
\end{enumerate}
\end{RiemannRemark}

\begin{RiemannExample}\label{ex:coverings-of-circle-torus-projective-space}
\begin{enumerate}[label=(\arabic*),ref=\arabic*,leftmargin=*]
\item\label{ex:real-line-covers-circle} 由
\[
  \pi(t)=(\cos(2\pi t),\sin(2\pi t))
\]
给出的映射 $\pi:\mathbb{R}\to S^1$ 是 $S^1$ 的一个光滑覆盖. 然而, 这个映射限制到 $(0,+\infty)$ 后是满射的局部微分同胚, 却不是覆盖映射.
\item\label{ex:euclidean-space-covers-torus} 覆盖映射的乘积显然是覆盖映射. 因此, 可以推广上例, 得到以 $\mathbb{R}^n$ 覆盖 $T^n\cong S^1\times\cdots\times S^1$ 的覆盖.
\item\label{ex:sphere-covers-projective-space} 在 Example~\ref{ex:projective-spaces-and-tori-as-quotients}(\ref{ex:projective-space-antipodal-quotient}) 中, 得到了以 $S^n$ 覆盖 $\mathbb{R}P^n$ 的覆盖.
\end{enumerate}
\end{RiemannExample}

设 $M$ 是覆盖流形. 若微分同胚 $h:M\to M$ 满足 $\pi\circ h=\pi$, 或等价地, $h$ 将每个集合 $\pi^{-1}(p)$ 映到自身, 则称 $h$ 为\textbf{甲板变换} (或\textbf{覆盖变换}). 可以证明, 全体覆盖变换所成的群 $G$ 是离散 Lie 群, 它在 $M$ 上的作用既自由又固有.

若覆盖流形 $M$ 是\textbf{单连通的} (参见第~\RGref{subsec:notes-proper-maps-and-fundamental-group}~节), 则称这个覆盖为\textbf{万有覆盖}. 此时, $B$ 与 $M/G$ 微分同胚. 此外, $G$ 同构于 $B$ 的\textbf{基本群} $\pi_1(B)$ (参见第~\RGref{subsec:notes-proper-maps-and-fundamental-group}~节).

\textbf{Lie 定理} (例如见 \RGcite{DK99}) 断言, 对给定的 Lie 代数 $\mathfrak g$, 存在唯一的单连通 Lie 群 $\widetilde G$, 其 Lie 代数为 $\mathfrak g$. 若 Lie 群 $G$ 的 Lie 代数也是 $\mathfrak g$, 则存在唯一的 Lie 群同态 $\pi:\widetilde G\to G$, 它是覆盖映射. 此时甲板变换群就是 $\ker(\pi)$, 因而 $G$ 与 $\widetilde G/\ker(\pi)$ 微分同胚. 实际上, $G$ 也同构于 $\widetilde G/\ker(\pi)$; 后者具有自然的群结构 (因为 $\ker(\pi)$ 是正规子群).

\begin{RiemannExample}\label{ex:fundamental-groups-from-universal-coverings}
\begin{enumerate}[label=(\arabic*),leftmargin=*]
\item 在 Example~\ref{ex:coverings-of-circle-torus-projective-space}(\ref{ex:real-line-covers-circle}) 的 $S^1$ 的万有覆盖中, 甲板变换是整数 $k$ 所给的平移 $h_k:t\mapsto t+k$, 因而 $S^1$ 的基本群是 $\mathbb{Z}$.
\item 类似地, $T^n$ 的万有覆盖的甲板变换是按整数向量的平移 [参见 Example~\ref{ex:coverings-of-circle-torus-projective-space}(\ref{ex:euclidean-space-covers-torus})], 因而 $T^n$ 的基本群是 $\mathbb{Z}^n$.
\item 在 Example~\ref{ex:coverings-of-circle-torus-projective-space}(\ref{ex:sphere-covers-projective-space}) 的 $\mathbb{R}P^n$ 的万有覆盖中, 仅有的甲板变换是恒等映射和对径映射, 因而 $\mathbb{R}P^n$ 的基本群是 $\mathbb{Z}_2$.
\end{enumerate}
\end{RiemannExample}
